% Section 5.7. Layer dilation identities and the associated coupled scale flow.
\subsection{Layer RG Beta Functions}
\label{sxn:sgd_layer_flow}

Our spectral RG separates a layer's overall amplitude from its
normalized spectrum. The same amplitude appears directly in
backpropagation: rescaling every weight in a layer turns its loss
response into a logarithmic scale derivative. These derivatives
define coupled dilation flows through one common loss.

Let $s_\ell=\ln c_\ell$ measure a layer dilation,
$\mathbf z_\ell$ its weighted input, and
$\mathcal L_{\mathcal B}=\langle L\rangle_{\mathcal B}$ the mean
prediction loss. Holding the minibatch, layer shapes, and biases
fixed, use this loss as a potential to define a continuous flow
in the log scales $\mathbf s=(s_1,\ldots,s_K)$, with parameter $\tau$:
\begin{tcolorbox}[marginalbox,title=The coupled dilation beta functions]
\eqtight{\begin{aligned}
 \beta_\ell(\mathbf s)
   :=\frac{ds_\ell}{d\tau}
   &=-\frac{\partial\mathcal L_{\mathcal B}}{\partial\ln c_\ell}\\
   &=-\left\langle\sum_i z_{\ell i}
               \frac{\partial L}{\partial z_{\ell i}}\right\rangle_{\mathcal B}\\
   &=-\langle\W_\ell,\nabla_{\W_\ell}\mathcal L_{\mathcal B}\rangle_F .
\end{aligned}}
\takeaway{Backpropagation supplies the logarithmic scale derivative
that generates each layer's loss-decreasing dilation.}
\end{tcolorbox}

\paragraph{Loss and layer scale.}
For a feedforward network, write $L^{(n)}$ for the prediction loss of
example $n$. Its minibatch mean is
\[
  \mathcal L_{\mathcal B}
   :=\frac{1}{|\mathcal B|}\sum_{n\in\mathcal B}L^{(n)}
   =\langle L\rangle_{\mathcal B},
\]
and its weight derivative is
$\nabla_{\W_\ell}\mathcal L_{\mathcal B}$.
Parameterize each layer by
\[
  \W_\ell=e^{s_\ell}\overline{\W}_\ell,\qquad
  c_\ell=e^{s_\ell}>0,\qquad
  \mathbf z_\ell=\W_\ell\mathbf h_{\ell-1} .
\]
Each fixed nonzero $\overline{\W}_\ell$ specifies the layer shape;
varying $s_\ell$ changes only its common amplitude.

\paragraph{The logarithmic derivative in backpropagation.}
For one example, suppress its index. The activation is
$\mathbf h_\ell=f_\ell(\mathbf z_\ell+\mathbf b_\ell)$.
Varying only $s_\ell$ gives
$\partial_{s_\ell}\mathbf z_\ell=\mathbf z_\ell$.
Where the network is differentiable, the chain rule gives
\[
\begin{aligned}
 \nabla_{\W_\ell}L
   &=(\nabla_{\mathbf z_\ell}L)\mathbf h_{\ell-1}^{\top},\\
 \frac{\partial L}{\partial s_\ell}
  &=\mathbf z_\ell^{\top}\nabla_{\mathbf z_\ell}L
   =\langle\W_\ell,\nabla_{\W_\ell}L\rangle_F .
\end{aligned}
\]
For each nonzero component, hold its sign and all other components
fixed. Its logarithmic derivative is
\[
 z_{\ell i}\frac{\partial L}{\partial z_{\ell i}}
   =\frac{\partial L}
          {\partial\ln|(\W_\ell\mathbf h_{\ell-1})_i|} .
\]
The response to a common scale of the weighted-input vector sums
these component responses. Averaging over examples gives
$\partial\mathcal L_{\mathcal B}/\partial\ln c_\ell$.
The operator $\partial_{\ln c_\ell}$ generates an infinitesimal
dilation, the rescaling part of RG. The boxed beta functions turn
this loss response into a flow along each layer's dilation coordinate.

\newpage
\paragraph{One loss couples the layer flows.}
Each $\beta_\ell(\mathbf s)$ responds to the other layer scales
through the common loss. The input $\mathbf h_{\ell-1}$ depends on
earlier layers, while
$\nabla_{\mathbf z_\ell}L$ carries the response through later layers.
Changing one layer can therefore change the responses of others.
A stationary scale for layer $\ell$ satisfies
$\beta_\ell(\mathbf s)=0$; its location can depend on the other
layers. A layer can be stationary while its neighbors still move,
and their motion can shift its stationary scale.

\begin{center}
\resizebox{\columnwidth}{!}{%
\begin{tikzpicture}[x=1cm,y=1cm,font=\sffamily\small,>=Stealth]
 \path[use as bounding box] (-.15,-.12) rectangle (8.05,3.75);
 \node[font=\sffamily\bfseries\small,text=boxblue]
   at (3.95,3.54) {One common minibatch loss};
 \node[font=\normalsize] (loss) at (3.95,3.08)
   {$\mathcal L_{\mathcal B}(s_1,s_2,s_3)$};
 \foreach \x/\l in {1.1/1,3.95/2,6.8/3}{
   \node[draw=boxblue,fill=boxlightblue,rounded corners=2pt,
     minimum width=2.2cm,minimum height=1.00cm,
     align=center,line width=.7pt] (layer\l) at (\x,2.03)
     {layer $\l$\\[3pt]$\dot s_{\l}=\beta_{\l}(\mathbf s)$};
   \draw[->,boxblue!65,line width=.65pt]
     (loss.south) -- (layer\l.north);
 }
 \draw[<->,boxblue,line width=1pt] (layer1.east) -- (layer2.west);
 \draw[<->,boxblue,line width=1pt] (layer2.east) -- (layer3.west);
 \draw[boxblue!55] (layer1.south) -- (1.1,.92) -- (6.8,.92) -- (layer3.south);
 \draw[boxblue!55] (layer2.south) -- (3.95,.92);
 \draw[->,boxblue!70] (3.95,.92) -- (3.95,.60);
 \node[font=\small] at (3.95,.30)
   {$\beta_1(\mathbf s^\star)=\beta_2(\mathbf s^\star)
     =\beta_3(\mathbf s^\star)=0$};
\end{tikzpicture}}
\captionof*{figure}{\textbf{Coupled layer dilations.}
Each layer contributes one scale response to the shared loss.
A joint stationary point satisfies every layer's dilation condition.}
\end{center}

The coupled flow decreases the common minibatch loss,
\[
 \frac{d\mathcal L_{\mathcal B}}{d\tau}
   =-\sum_{\ell=1}^{K}\beta_\ell(\mathbf s)^2\leq0 .
\]
At a joint fixed point, $\boldsymbol\beta(\mathbf s^\star)=0$, so
no layer dilation changes the loss to first order. The loss itself
need not vanish.

\paragraph{The same dilation rescales the ESD.}
For $\W_\ell\in\mathbb R^{N_\ell\times M_\ell}$, the correlation matrix is
$\X_\ell=\W_\ell^\top\W_\ell/N_\ell$.
A dilation acts on this operator, its eigenvalues, and its
probability-normalized ESD as
\[
\begin{aligned}
 \W_\ell&\mapsto c_\ell\W_\ell,&
 \X_\ell&\mapsto c_\ell^2\X_\ell,\\
 \lambda_{\ell i}&\mapsto c_\ell^2\lambda_{\ell i},&
 \rho_{\ell,c}(\lambda)&=c_\ell^{-2}
                 \rho_\ell(\lambda/c_\ell^2).
\end{aligned}
\]
Along the flow, $\partial_{s_\ell}\X_\ell=2\X_\ell$ gives
\[
 \frac{d\X_\ell}{d\tau}=2\beta_\ell(\mathbf s)\X_\ell,\qquad
 \frac{d\ln\lambda_{\ell i}}{d\tau}=2\beta_\ell(\mathbf s)
\]
for $\lambda_{\ell i}>0$ in the second identity. Thus every
positive eigenvalue shifts at the same rate on a logarithmic axis.
The full-spectrum first-moment normalization of
Section~\ref{sxn:meanlog_gauge} removes this common amplitude,
\[
 \widetilde{\X}_\ell=\frac{\X_\ell}{m_\ell},\qquad
 m_\ell=\frac{1}{M_\ell}\Tr\X_\ell,\qquad
 \partial_{\ln c_\ell}\widetilde{\X}_\ell=0 .
\]
Both $\X_\ell$ and its positive mean $m_\ell$ acquire the factor
$c_\ell^2$, so their ratio is unchanged. The normalized ESD retains
the spectral shape used for ECS selection. Thus backpropagation
exposes the same layer dilation whose amplitude our spectral RG
normalization removes.

